% !TeX root = ../se320_structural_assurability.tex
% ==================================================
\section{Introduction}
\label{aa.sec.intro}
% ==================================================

Assurance depends on evidence.

For a claim that a system is safe, secure, reliable, compliant, or otherwise trustworthy,
an evaluator must be able to obtain evidence bearing on that claim and connect the evidence to the claim through a defensible argument.
Assurance-case methods make this relationship explicit: claims are supported by evidence through structured reasoning,
while separate questions concern the adequacy of the argument and the confidence warranted
by it~\cite{goodenough2012confidence,bloomfield2010safetycases}.

For increasingly capable AI systems, however, an important question arises before the assurance argument can even be constructed:

\begin{quote}
  What evidence does the structure of the system make possible to obtain?
\end{quote}

Two AI systems may exhibit equivalent externally measured performance while differing substantially
in their capacity to expose evidence concerning how that performance was produced.
One system may expose only prompts and outputs.
Another may additionally preserve tool invocations, retrieval events, authorization decisions,
state transitions, provenance, independently generated records, and selected internal-state measurements.
Even when the systems behave identically over a specified evaluation suite, they need not support the same assurance inquiries.

This distinction becomes increasingly important as AI systems move from isolated prediction toward persistent,
tool-using, stateful, and agentic operation.
Pre-deployment evaluation remains valuable, but deployed systems operate under changing inputs, environments, integrations, and incentives.
Recent work has accordingly expanded from benchmark evaluation toward runtime monitoring,
agent observability, auditability, provenance, assurance cases, interpretability, and
internal-state monitoring \cite{nist2026monitoring,fernsel2024auditability,openai2025cotmonitorability,nistAgenticProbes2026}.
These efforts expose a common underlying problem: assurance depends not merely on the existence of evaluation methods,
but on whether the system and its operating environment make the evidence required by those methods obtainable.

This problem has mature analogues elsewhere in engineering.

Hardware design for testability recognizes that testing effectiveness is partly determined by the structure of the device being tested.
Observability and controllability can be deliberately designed into hardware so that relevant internal states
can be established and examined economically \cite{abramovici1990digitaltesting,williams1983testability}.
Software architecture similarly affects testability, analyzability, diagnosability, and the practical cost of
verification \cite{bass2021softwarearchitecture,clements2002evaluating,kazman2000atam}.
Safety engineering recognizes that architecture affects the generation and independence of assurance evidence,
while security engineering emphasizes provenance, tamper resistance, privilege separation, and evidence integrity
under adversarial conditions \cite{saeARP4754A,rtcaDO178C,saltzer1975protection,anderson2020security}.

AI assurance draws on all of these traditions, but the corresponding structural property is not usually the primary object of analysis.

A system may be described as auditable, monitorable, interpretable, observable, testable, or capable of supporting an assurance case.
Each concept captures an important part of the problem.
None alone answers the more general comparative question considered here:

\begin{quote}
  Given two AI systems subject to the same assurance claim, what structural differences determine
  whether one permits assurance-relevant evidence to be obtained that the other does not?
\end{quote}

We call this property \emph{structural assurability}.

Structural assurability is the claim-relative capacity of a system and its evidence-producing environment
to make evidence relevant to supporting or refuting an assurance claim practically obtainable
under specified access, cooperation, and resource constraints.

We write this dependence schematically as

\begin{equation}
  \mathcal{A}(S,C,K,B),
\end{equation}

where \(S\) is the system, \(C\) is the assurance claim, \(K\) represents applicable knowledge,
access, trust, and cooperation conditions,
and \(B\) represents practical resource constraints.

The formulation is intentionally claim-relative.

A system that exposes extensive internal state may be highly assurable for a claim concerning internal computation while
gaining little advantage for a claim resolved completely from an externally observable output.
Likewise, an architecture with strong independent external monitoring may support some claims better
than an architecture offering richer but operator-controlled internal telemetry.

Accordingly,

\begin{equation}
  \text{assurability}(S)
\end{equation}

without a specified assurance context is generally underspecified.

Structural assurability is also deliberately distinct from safety or other favorable system properties.

A system can be structurally highly assurable and nevertheless fail the assurance claim under investigation.
Richer evidence may make failure easier to establish.
Structural assurability concerns capacity for examination rather than the result of examination:

\begin{equation}
  \boxed{
    \begin{aligned}
      \text{Property satisfaction}
       & \neq
      \text{Structural assurability} \\
       & \neq
      \text{Assurance argument}
      \neq
      \text{Confidence in the argument}.
    \end{aligned}
  }
\end{equation}

This separation is important because otherwise evidence-producing mechanisms can be conflated with
evidence that happens to support a desired conclusion.
A useful assurance structure should support both confirmation and refutation.

The central premise of this paper is therefore modest:

\begin{quote}
  Systems with equivalent functional behavior or measured performance can differ materially in assurability because their
  architectures and evidence-producing environments can differ in their capacity to make assurance-relevant distinctions possible.
\end{quote}

The phrase \emph{assurance-relevant distinctions} provides the core interpretation developed in this work.

Consider a claim concerning whether an AI agent performed an unauthorized consequential action.
Relevant alternatives may include no action, authorized agent action, unauthorized agent action,
human-initiated action, externally initiated action, or an action whose historical record has been altered.
An evidence-producing structure contributes to assurability
when it makes one or more such alternatives distinguishable where they would otherwise remain unresolved.

This formulation avoids equating assurability with evidence volume.

More logs do not necessarily imply greater assurability.
More internal state does not necessarily imply greater assurability.
More evaluation access does not necessarily imply greater assurability.
A structural difference matters when it enables a distinction material to the assurance claim.

We initially identify seven structural dimensions that can affect this capacity:

\begin{equation}
  \begin{aligned}
    D_S = \{ &
    \text{observability},
    \text{controllability},
    \text{traceability},
    \text{reconstructability}, \\
             &
    \text{independence},
    \text{integrity},
    \text{evidence coverage}
    \}.
  \end{aligned}
\end{equation}

These dimensions are provisional.
We do not claim that they are mutually independent,
universally required, jointly sufficient, or complete.

Instead, the paper adopts a weaker criterion.

A structural dimension is assurance-relevant when there exists an assurance context
and a pair of otherwise suitably controlled systems for which variation in that dimension creates a material difference
in the assurance-relevant distinctions supported by obtainable evidence.

We refer to such a pair as a \emph{separating construction}.

The resulting logical claim is therefore not

\begin{equation}
  \text{all seven dimensions}
  \Rightarrow
  \text{sufficient assurance}.
\end{equation}

Rather, the criterion for establishing a candidate dimension \(d\)
requires a specified assurance context and a controlled pair
of systems exhibiting both a difference in the dimension
and an independently justified difference in claim-material
evidentiary capabilities.

In this sense, omitting \(d\) from a general model can collapse systems that differ materially in assurability for at least one assurance context.

The framework is consequently open-ended.
If a future architecture, threat model, assurance claim, or measurement mechanism reveals an
additional structural property that admits a separating construction not representable by the current dimensions,
that property can be added without invalidating the existing framework.

The paper also avoids defining a universal scalar assurability score.

Different systems may expose different assurance-relevant capabilities without either capability set containing the other.
An architecture with strong internal-state observability but weak evidence independence may be incomparable
to an architecture with limited internal visibility but independently generated, tamper-resistant records of consequential external actions.

The appropriate comparative structure is therefore generally a partial relation admitting
equivalence, dominance, and incomparability rather than a total ranking.

This paper makes the following contributions:

\begin{enumerate}


  \item It distinguishes \emph{structural assurability} from property satisfaction,
        assurance arguments, confidence in assurance arguments, auditability, monitorability, and related concepts.

  \item It defines assurability as a claim-relative property under explicit access, cooperation, and resource constraints.

  \item It separates structural properties of an evidence-producing architecture from downstream
        evidence quality, measurement validity, uncertainty, and assurance reasoning.

  \item It identifies an initial set of seven candidate structural dimensions:
        observability, controllability, traceability, reconstructability, independence, integrity, and evidence coverage.

  \item It introduces a separating-construction criterion for demonstrating that a proposed structural dimension
        contributes to assurability without claiming that the dimension is universally necessary or that the resulting dimension set is sufficient.

  \item It develops a comparative model in which systems are evaluated according to the assurance-relevant distinctions
        made possible by practically obtainable evidence, allowing equivalence, strict dominance, and incomparability without requiring a universal scalar score.

  \item It applies the framework to representative AI architectures and demonstrates that richer instrumentation
        can increase assurability for one claim, provide no additional assurability for another, or create tradeoffs that prevent a total ordering.


\end{enumerate}

The constituent ideas developed here are not individually new.
Observability, controllability, testability, auditability, monitorability, assurance cases,
provenance, evidence integrity, independence, and evaluation rigor each have substantial intellectual histories.

The proposed contribution is their integration around a different object of study:
the structural capacity of an AI system and its evidence-producing environment to make claim-relevant distinctions possible.

The remainder of the paper proceeds as follows.
Section~\ref{aa.sec.prior-work} examines the foundations of the framework
across assurance theory, design for testability, software architecture, safety engineering,
security evaluation, auditability, AI monitoring, interpretability, measurement science, and related work.
Section~\ref{aa.sec.foundations} defines structural assurability and develops the distinction among structural properties,
evidence properties, and assurance constraints.
Section~\ref{aa.sec.formal-model} formalizes comparative assurability and establishes the separating-construction criterion.
Section~\ref{aa.sec.applications} applies the framework to representative AI architectures.
Section~\ref{aa.sec.theoretical-results} records what is established in the Lean formalization.
Section~\ref{aa.sec.pilot} describes an illustrative application to an external evaluation framework.
Sections~\ref{aa.sec.discussion} and~\ref{aa.sec.threats} discuss implications, limitations, and threats to validity.
Section~\ref{aa.sec.conclusion} concludes.
